\documentclass[11pt]{article}
\usepackage[margin=1in]{geometry}
\usepackage{amsmath,amssymb,amsthm}
\usepackage{hyperref}
\usepackage[T1]{fontenc}
\usepackage[utf8]{inputenc}

\newtheorem{definition}{Definition}
\newtheorem{principle}{Design Principle}
\newtheorem{milestone}{Milestone}
\newtheorem{risk}{Risk}

\title{Medium-Scale PeTTa Memory for ProtomegaTron:\\
Detailed Design and Step-by-Step Implementation Plan}
\author{OmegaClaw / ZeroBot working notes for Ben Goertzel}
\date{2026-06-27}

\begin{document}
\maketitle

\begin{abstract}
This ASCII-only LaTeX note specifies a conservative design and implementation plan
for adding an intermediate-scale PeTTa memory layer to ProtomegaTron.  The layer
sits between volatile working memory and broad ChromaDB long-term memory.  Its
purpose is to preserve structured, queryable, provenance-bearing cognitive state:
episodes, entities, commitments, hypotheses, open questions, action boundaries,
observations, and distillation links.  The plan is intentionally incremental: first
add a file-backed atom store and read-only query surface, then add selective
extraction, then add distillation and eventual integration into the OmegaClaw loop.
The 2026-06-27 revision makes the design explicitly PLN-ready: the same medium
memory should later support inference over a normalized, truth-valued AtomSpace
view without confusing quoted conversation, observed events, decisions, hypotheses,
and derived beliefs.
\end{abstract}

\section{Motivation}

\subsection{Revision summary}

This version folds in three constraints that should shape the implementation:
\begin{enumerate}
  \item The canonical unit is a \texttt{MemoryCluster}, not an isolated atom line.
  \item Observed events, quoted utterances, commitments, hypotheses, decisions, and
        derived beliefs must remain epistemically distinct.
  \item The store is append-only; current state is a derived view over status,
        supersession, and truth-value events.
\end{enumerate}

The current OmegaClaw memory architecture, as inspected in the local
ProtomegaTron/OmegaClaw checkout, has four important memory surfaces:
\begin{enumerate}
  \item Working memory, represented by \texttt{pin}, useful for immediate task state.
  \item The episodic text trace, \texttt{memory/history.metta}, whose tail is injected
        into prompts and whose neighborhoods can be retrieved by \texttt{episodes}.
  \item ChromaDB long-term memory, accessed by \texttt{remember} and \texttt{query},
        useful for semantic recall across longer time spans.
  \item Per-invocation MeTTa/AtomSpace reasoning, useful but not persistent across
        separate calls unless the agent manually carries atoms forward.
\end{enumerate}

This leaves a gap.  ProtomegaTron needs a memory layer that is more precise than
embedding recall, more persistent than a single pin or one fresh AtomSpace call,
and more structured than a raw text history tail.  The proposed layer is a
medium-scale PeTTa memory, abbreviated here as \texttt{MM}.

\section{Design goal}

\begin{definition}[Medium-scale PeTTa memory]
A medium-scale PeTTa memory is a bounded, persistent, typed atom store
\[
  MM_t = \{ a_1, a_2, \ldots, a_n \}
\]
where each atom or small atom cluster records a currently useful cognitive item:
an episode, entity, commitment, hypothesis, question, boundary, observation,
artifact, or provenance relation.  It is intended for active structured cognition,
not for permanent archival storage of everything the agent sees.
\end{definition}

\begin{principle}[Selective atomization]
Only high-value events should enter \texttt{MM}.  The store should contain facts
that are likely to affect near-future reasoning, task continuity, self-scrutiny,
or Hyperseed formalization.  It should not become a complete transcript mirror.
\end{principle}

\begin{principle}[Provenance first]
Every nontrivial assertion should either include direct provenance or link to a
provenance atom.  A false or stale symbolic assertion without provenance is more
dangerous than an unstructured text memory, because it looks more formal than it is.
\end{principle}

\begin{principle}[Epistemic role separation]
Atoms that record what happened should be separated from atoms that claim what is
true.  A message from Ben, a quote from an LLM, a derived belief, a standing
decision, and a speculative hypothesis should have different predicates and should
not be collapsed into a single untyped fact.  This is especially important once PLN
inference is run over the memory AtomSpace: quoted claims may provide evidence, but
they should not automatically become premises.
\end{principle}

\begin{principle}[Append-only audit semantics]
The canonical store should append new status, salience, and truth-value events
rather than silently mutating old atoms.  Current status can be derived from the
latest non-superseded status event.  This preserves an audit trail and makes
contradictions inspectable instead of hidden.
\end{principle}

\begin{principle}[Bounded autonomy]
The first implementation should not grant ProtomegaTron new external authority.  It
should only write a local structured memory file and read/query that file.  Telegram,
GitHub, shell, paid compute, and filesystem permissions remain governed by existing
boundaries.
\end{principle}

\begin{principle}[PLN-ready AtomSpace]
Ben noted on 2026-06-27 that later versions should run PLN inference on this memory
AtomSpace.  Therefore version 1 should avoid ad hoc string-only structures that would
be hard to promote into PLN premises.  Atoms should use stable predicates, explicit
truth values for interpretive claims when appropriate, provenance links, and status
events that can later be loaded into a reasoning context without lossy translation.
\end{principle}

\section{Placement in the architecture}

The intended memory stack is:
\[
  WM \rightarrow MM \rightarrow LTM
\]
where:
\begin{itemize}
  \item \texttt{WM} is working memory: pins, current loop state, current prompt tail.
  \item \texttt{MM} is medium PeTTa memory: structured atoms over recent and active
        concerns.
  \item \texttt{LTM} is ChromaDB plus durable project records and curated Markdown.
\end{itemize}

Operationally, a turn may use the layers as follows:
\begin{enumerate}
  \item Receive message or autonomous wake event.
  \item Retrieve relevant \texttt{MM} atoms by type, entity, status, or recency.
  \item Retrieve relevant ChromaDB memories if broad semantic context is needed.
  \item Construct prompt context from mandate, skills, selected history, selected
        \texttt{MM} atoms, and recent results.
  \item Let the LLM propose skill calls and optional memory updates.
  \item Execute safe skill calls.
  \item Append high-value structured atoms to \texttt{MM}.
  \item Periodically distill older \texttt{MM} atoms into ChromaDB or project notes.
\end{enumerate}

\subsection{Canonical log and derived views}

Version 1 should distinguish four products:
\begin{enumerate}
  \item \texttt{MM-log}: append-only canonical atom journal.  This is the source of
        truth and should be small enough to inspect with ordinary tools.
  \item \texttt{MM-index}: rebuildable query acceleration, such as type, id, about,
        status, and time indexes.  It can be deleted and regenerated.
  \item \texttt{MM-prompt-view}: bounded, human-readable context injected into LLM
        prompts.  It should never be treated as canonical.
  \item \texttt{MM-pln-view}: normalized, typed, truth-valued atoms safe to load as
        PLN premises or evidence.  It should exclude raw quoted claims unless a
        rule explicitly promotes them into evidence-bearing beliefs.
\end{enumerate}

This separation lets the same substrate support auditability, ordinary prompt
continuity, and later probabilistic reasoning without requiring all consumers to
share the same representation.

\section{Minimal atom schema}

The schema should be small enough to remember and inspect by eye.  The first version
should use flat, readable predicates rather than deep ontological machinery.

\subsection{Identifiers}

Use stable, opaque identifiers:
\begin{verbatim}
(Episode e-20260627-121018-001)
(Entity ent-protomegatron)
(Question q-medium-memory-decay)
(Hypothesis h-medium-memory-useful)
(Commitment c-write-medium-memory-plan)
(Artifact art-medium-memory-plan-tex)
(Boundary b-telegram-allowlist)
(MemoryCluster mc-20260627-121018-001)
(Source src-loop-metta-16d380d)
\end{verbatim}

Identifier construction rule:
\begin{verbatim}
<kind>-<yyyymmdd>-<hhmmss>-<counter>
\end{verbatim}
for event-like records, and stable slugs for durable entities or code anchors.

\subsection{Cluster envelope}

The canonical write unit should be a small cluster, not an isolated line.  This makes
append-only updates, provenance inspection, and later PLN loading safer.

\begin{verbatim}
(MemoryCluster mc1)
(ClusterType mc1 episode-record)
(ClusterOpenedAt mc1 "2026-06-27 12:10 PDT")
(ClusterSource mc1 src-telegram-history-20260627)
(Contains mc1 e1)
(Contains mc1 c1)
(Contains mc1 q1)
(ClusterStatus mc1 active)
\end{verbatim}

Cluster conventions:
\begin{enumerate}
  \item Every append writes one complete cluster.
  \item A later correction writes a new cluster and links it with \texttt{Supersedes},
        rather than editing the old cluster in place.
  \item Queries may return whole clusters even when only one contained atom matched.
  \item The PLN view may select only premise-safe atoms from a cluster.
\end{enumerate}

\subsection{Epistemic roles}

The first schema should reserve explicit predicates for epistemic roles, and should
tag clusters or atoms when the role is not obvious from the predicate name.  This
makes the memory useful for ordinary retrieval now and safer for PLN inference later.

\begin{verbatim}
(ObservedEvent e1)
(SpeechEvent se1)
(QuotedClaim qc1)
(DerivedBelief b1)
(Decision d1)
(Hypothesis h1)
(OpenQuestion q1)
(Commitment c1)
\end{verbatim}

The intended discipline is:
\begin{itemize}
  \item \texttt{ObservedEvent}: something the local system directly observed, such as
        a Telegram message, file creation, test result, or tool output.
  \item \texttt{SpeechEvent}: a communicative event, usually linking actor, channel,
        and message id.
  \item \texttt{QuotedClaim}: what an actor or model asserted.  It is evidence, not
        automatically a fact.
  \item \texttt{DerivedBelief}: a normalized proposition that the agent is willing to
        use as a belief, with a truth value and evidence links.
  \item \texttt{Decision}: an adopted course of action with rationale and provenance.
  \item \texttt{Hypothesis}: a plausible but unsettled proposition, normally with a
        lower confidence than a decision or directly observed fact.
\end{itemize}

\subsection{Episode atoms}

\begin{verbatim}
(Episode e1)
(AtTime e1 "2026-06-27 12:10 PDT")
(Channel e1 telegram-group)
(Actor e1 Ben)
(ObservedBy e1 ProtomegaTron)
(Said e1 Ben "Great please progress w this...")
(About e1 ProtomegaTronDesign)
(About e1 MediumPeTTaMemory)
(HasStatus e1 recorded)
(HasProvenance e1 src-telegram-history-20260627)
(EpistemicRole e1 observed-event)
\end{verbatim}

For PLN-readiness, prefer the more explicit event-plus-claim form when the content
of a message may later be interpreted as evidence:

\begin{verbatim}
(ObservedEvent e1)
(SpeechEvent se1)
(PartOf se1 e1)
(AtTime se1 "2026-06-27 14:02 PDT")
(Channel se1 telegram-group)
(Speaker se1 Ben)
(MessageId se1 telegram-msg-505)
(RawUtterance se1 "Note that later we will be running PLN inference...")
(QuotedClaim qc1)
(ClaimSource qc1 se1)
(ClaimText qc1 "Later versions should run PLN inference on this memory AtomSpace")
(About qc1 MediumPeTTaMemory)
(About qc1 PLN)
(ClaimStatus qc1 quoted-unverified)
\end{verbatim}

\subsection{Commitment atoms}

\begin{verbatim}
(Commitment c1)
(CommittedBy c1 ProtomegaTron)
(CommittedTo c1 Ben)
(CommitmentText c1 "Write an ASCII LaTeX design and implementation plan")
(CommitmentStatus c1 active)
(OpenedBy c1 e1)
(HasDueMode c1 immediate-next-step)
(HasProvenance c1 e1)
(EpistemicRole c1 commitment)
\end{verbatim}

When status changes, append status events instead of overwriting the old status:

\begin{verbatim}
(StatusEvent st1)
(StatusSubject st1 c1)
(StatusValue st1 active)
(StatusTime st1 "2026-06-27 13:51 PDT")
(StatusProvenance st1 e1)

(StatusEvent st2)
(StatusSubject st2 c1)
(StatusValue st2 resolved)
(StatusTime st2 "2026-06-27 14:10 PDT")
(Supersedes st2 st1)
(StatusProvenance st2 art-medium-memory-plan-tex)
\end{verbatim}

\subsection{Question atoms}

\begin{verbatim}
(OpenQuestion q1)
(QuestionText q1 "How should medium PeTTa memory decay and distill?")
(QuestionStatus q1 active)
(QuestionScope q1 architecture)
(About q1 MediumPeTTaMemory)
(HasProvenance q1 e1)
(EpistemicRole q1 open-question)
\end{verbatim}

\subsection{Hypothesis atoms}

Truth values can be represented as ordinary MeTTa structures initially, without
requiring every query to invoke PLN.

\begin{verbatim}
(Hypothesis h1)
(HypothesisText h1 "Selective medium PeTTa memory improves continuity")
(About h1 ProtomegaTronMemory)
(Confidence h1 (stv 0.80 0.60))
(EvidenceFor h1 e1)
(HypothesisStatus h1 active)
(EpistemicRole h1 hypothesis)
\end{verbatim}

For PLN, the proposition should be explicit rather than trapped inside a string:

\begin{verbatim}
(Hypothesis h1)
(HypothesisContent h1 (Improves SelectiveMediumPeTTaMemory AgentContinuity))
(TruthValue h1 (stv 0.80 0.60))
(EvidenceFor h1 qc1)
(HasProvenance h1 se1)
(StatusSubject st-h1-1 h1)
(StatusValue st-h1-1 active)
\end{verbatim}

The string form may still be kept for prompt readability, but the symbolic content
is what should be loaded into the PLN view.

\subsection{Belief atoms for PLN}

Derived beliefs are the main bridge between audit memory and inference memory:

\begin{verbatim}
(DerivedBelief b1)
(BeliefContent b1 (Requires MediumPeTTaMemory ExplicitTruthValues))
(TruthValue b1 (stv 0.90 0.70))
(EvidenceFor b1 qc1)
(EvidenceFor b1 design-review-20260627)
(BeliefStatus b1 active)
(HasProvenance b1 se1)
\end{verbatim}

Rules may later promote quoted claims into evidence-bearing beliefs when a source is
trusted in a relevant context, for example:

\begin{verbatim}
(TrustedSource Ben medium-memory-design)
(CanPromoteToEvidence qc1 b1)
\end{verbatim}

but promotion should be explicit and logged, not implicit.

\subsection{Salience and retention atoms}

Bounded memory requires explicit attention and retention metadata:

\begin{verbatim}
(SalienceEvent sal1)
(SalienceSubject sal1 c1)
(SalienceValue sal1 0.82)
(SalienceReason sal1 active-task)
(RetentionClass c1 active-task)
(ReviewAfter c1 "2026-06-28")
\end{verbatim}

The prompt view should prefer high-salience active items.  The distillation pass
should prefer low-salience resolved items and stale open questions.

\subsection{Boundary atoms}

\begin{verbatim}
(Boundary b1)
(BoundaryName b1 "Telegram allowlist")
(Constrains b1 send-action)
(Constrains b1 receive-action)
(ImplementedIn b1 "channels/telegram.py")
(HasStatus b1 active)
(HasProvenance b1 src-telegram-py-16d380d-local)
(EpistemicRole b1 action-boundary)
\end{verbatim}

\subsection{Artifact atoms}

\begin{verbatim}
(Artifact art1)
(ArtifactPath art1 "papers/0003-medium-petta-memory-plan/medium_petta_memory_plan.tex")
(ArtifactType art1 latex-source)
(About art1 MediumPeTTaMemory)
(CreatedBy art1 ProtomegaTron)
(HasProvenance art1 c1)
(EpistemicRole art1 artifact-record)
\end{verbatim}

\subsection{Append-only update atoms}

Status changes, confidence revisions, and corrections should be new atoms or new
clusters.  The current view is computed from the latest non-superseded update.

\begin{verbatim}
(StatusEvent se1)
(StatusSubject se1 c1)
(StatusValue se1 resolved)
(AtTime se1 "2026-06-27 12:24 PDT")
(HasProvenance se1 art1)

(TruthValueEvent tv1)
(TruthSubject tv1 h1)
(TruthValue tv1 (stv 0.85 0.65))
(Supersedes tv1 tv0)
(HasProvenance tv1 e2)
\end{verbatim}

This pattern is more verbose than mutation, but it preserves the audit trail needed
for both scientific provenance and later PLN belief revision.

\section{File layout}

The first prototype should avoid database complexity.  Use append-only files under the
OmegaClaw project or a local experimental branch:

\begin{verbatim}
OmegaClaw-Core/
  memory/
    medium_memory.metta          # append-only canonical atom log
    medium_memory_index.metta    # optional derived index, rebuildable
  src/
    medium_memory.metta          # MeTTa API wrapper
    medium_memory.py             # Python helper for simple file operations
\end{verbatim}

If implemented first in the research workspace rather than upstream OmegaClaw Core,
use:

\begin{verbatim}
projects/omegaclaw/experiments/medium-memory-prototype/
  medium_memory.metta
  medium_memory.py
  tests/
\end{verbatim}

\section{API design}

The first API should be deliberately small.

\subsection{Write API}

\begin{verbatim}
(mm-append $atom)
(mm-append-cluster $cluster)
(mm-record-episode $episode-id $cluster)
(mm-append-status $subject-id $status $provenance)
\end{verbatim}

Expected behavior:
\begin{enumerate}
  \item Validate that the atom is parseable as MeTTa syntax.
  \item Reject atoms above a configured size limit.
  \item Append with a timestamp comment or explicit \texttt{AtTime} atom.
  \item Return a success/failure status that can be fed back into the loop.
\end{enumerate}

\subsection{Read API}

\begin{verbatim}
(mm-tail $max-chars)
(mm-query-cluster $cluster-id)
(mm-query-type $type $limit)
(mm-query-about $entity $limit)
(mm-query-status $status $limit)
(mm-query-role $epistemic-role $limit)
(mm-query-id $id)
\end{verbatim}

The first implementation can use simple text scanning.  Later versions can load the
file into an AtomSpace or maintain derived indexes.

\subsection{View API}

The store should expose separate views for separate consumers:

\begin{verbatim}
(mm-prompt-view $topic $max-chars)
(mm-pln-view $topic $max-atoms)
(mm-audit-view $id $max-chars)
\end{verbatim}

Expected behavior:
\begin{itemize}
  \item \texttt{mm-prompt-view} returns a compact textual summary of active
        commitments, questions, salient observations, and relevant boundaries.
  \item \texttt{mm-pln-view} returns normalized typed atoms with explicit truth values
        and evidence links.  It should not include raw quoted claims as factual
        premises unless they have been explicitly promoted.
  \item \texttt{mm-audit-view} returns the provenance chain and append-only event
        history for one id.
\end{itemize}

This division is a safety feature, not merely an implementation detail: prompt
selection, human audit, and PLN inference have different failure modes.

\subsection{Distillation API}

\begin{verbatim}
(mm-candidates-for-distill $older-than $limit)
(mm-mark-distilled $id $destination)
(mm-export-summary $ids)
\end{verbatim}

Distillation should not delete records initially.  It should mark them as distilled
and link to the destination: ChromaDB memory id, project Markdown file, Git commit,
or formalization note.

\section{Prompt integration}

The prompt should not include the full \texttt{MM} store.  It should include a bounded
selection:
\begin{enumerate}
  \item Active commitments involving Ben or the current project.
  \item Active open questions about the current project or agent design.
  \item Recent episodes relevant to the current message.
  \item Boundaries relevant to proposed actions.
  \item Recent failed or surprising observations.
\end{enumerate}

A possible prompt block:
\begin{verbatim}
MEDIUM_MEMORY_ACTIVE:
- Commitment c1: write ASCII LaTeX design plan. Status active.
- Question q1: how should medium PeTTa memory decay and distill? Status active.
- Boundary b1: Telegram allowlist constrains send/receive actions. Status active.
\end{verbatim}

This human-readable summary can be generated from atoms.  The atoms remain the source
of truth; the summary is a prompt-facing view.

\section{Step-by-step implementation plan}

\begin{milestone}[Step 0: freeze scope]
Goal: implement a local experimental medium memory that cannot affect external actions
except by improving prompt context.

Deliverables:
\begin{itemize}
  \item A short design note, this file.
  \item A minimal schema file with example atoms.
  \item A decision on prototype location: local experiment first, then OmegaClaw Core.
\end{itemize}

Exit criteria:
\begin{itemize}
  \item No new credentials, services, or external write authority are introduced.
  \item The atom schema fits on one page and has examples.
\end{itemize}
\end{milestone}

\begin{milestone}[Step 1: create append-only store]
Implement \texttt{memory/medium\_memory.metta} and helper functions for appending
validated clusters and atoms.

Tasks:
\begin{enumerate}
  \item Add \texttt{src/medium\_memory.py} with \texttt{append\_atom},
        \texttt{append\_cluster}, \texttt{tail}, and parse-check helpers.
  \item Add \texttt{src/medium\_memory.metta} wrappers around those Python calls.
  \item Add limits: maximum atom length, maximum cluster length, maximum tail chars.
  \item Add cluster-envelope checks: every cluster needs a \texttt{MemoryCluster} id,
        \texttt{ClusterType}, timestamp, provenance/source, and at least one
        \texttt{Contains} link.
  \item Add tests for successful cluster append, malformed atom rejection, missing
        cluster metadata, and size rejection.
\end{enumerate}

Exit criteria:
\begin{itemize}
  \item A test can append an episode cluster and read it back.
  \item Malformed syntax does not corrupt the store.
  \item Existing OmegaClaw tests still pass or the prototype tests pass in isolation.
\end{itemize}
\end{milestone}

\begin{milestone}[Step 2: add simple query functions]
Implement text-based query functions before attempting full AtomSpace indexing.

Tasks:
\begin{enumerate}
  \item Query by id: return all lines mentioning the id.
  \item Query by cluster id: return the complete cluster.
  \item Query by type: return clusters beginning with \texttt{(Episode ...)},
        \texttt{(Commitment ...)}, \texttt{(OpenQuestion ...)}, etc.
  \item Query by \texttt{About}: return clusters containing \texttt{(About <id> <entity>)}.
  \item Query by status: active, blocked, resolved, stale, distilled.
  \item Query by epistemic role: observed-event, quoted-utterance, hypothesis,
        decision, action-boundary, derived-belief.
  \item Add bounded output limits for prompt safety.
\end{enumerate}

Exit criteria:
\begin{itemize}
  \item The agent can retrieve active commitments and active questions.
  \item Queries are deterministic and bounded.
  \item Query failure returns explicit error feedback, not an empty hallucination trap.
\end{itemize}
\end{milestone}

\begin{milestone}[Step 3: prompt-view selection]
Generate a compact prompt-facing medium-memory view.

Tasks:
\begin{enumerate}
  \item Add \texttt{getMediumMemoryContext} in MeTTa or Python.
  \item Select active commitments, active questions, relevant boundaries, and recent
        high-value episodes.
  \item Insert the resulting bounded text into \texttt{getContext} after history or
        before history.
  \item Add a config knob: \texttt{mediumMemoryMaxContextChars}.
  \item Default the feature off until tests pass; then enable for ProtomegaTron only.
\end{enumerate}

Exit criteria:
\begin{itemize}
  \item A controlled mock run shows the prompt contains selected \texttt{MM} facts.
  \item Idle/no-input loops still do not trigger backend calls.
  \item Prompt size remains bounded and observable in logs.
\end{itemize}
\end{milestone}

\begin{milestone}[Step 4: selective extraction protocol]
Teach ProtomegaTron when to write medium-memory atoms.

Tasks:
\begin{enumerate}
  \item Add prompt guidance: write \texttt{MM} atoms only for commitments, open
        questions, decisions, hypotheses, important observations, and boundaries.
  \item Add prompt guidance: separate observed events and quoted claims from derived
        beliefs; do not write a conversational quote as a naked fact.
  \item Add a skill such as \texttt{medium-remember} that appends a validated cluster.
  \item Create examples for common cases: Ben request, completed artifact, blocker,
        design hypothesis, decision, stale belief correction.
  \item Require explicit epistemic roles for quoted utterances, hypotheses, decisions,
        and derived beliefs.
  \item Add tests using the mock LLM or deterministic helper outputs.
\end{enumerate}

Exit criteria:
\begin{itemize}
  \item A fresh Ben request creates an episode and commitment cluster.
  \item A design assertion creates a speech event and quoted claim; an optional
        derived belief is added only with truth value and evidence links.
  \item A completion update marks the commitment resolved and links an artifact.
  \item Routine greetings do not create medium-memory atoms unless they establish
        a durable fact or task.
\end{itemize}
\end{milestone}

\begin{milestone}[Step 5: distillation]
Add a non-destructive distillation pass.

Tasks:
\begin{enumerate}
  \item Select old resolved commitments and stale questions as candidates.
  \item Summarize candidates into a ChromaDB \texttt{remember} string or project
        Markdown entry, with provenance links.
  \item Mark source atoms as distilled but do not delete them.
  \item Add a report command showing what would be distilled before writing.
\end{enumerate}

Exit criteria:
\begin{itemize}
  \item Distillation can be run manually and reviewed.
  \item The ChromaDB memory or project record links back to source atom ids.
  \item No atom is deleted in version 1.
\end{itemize}
\end{milestone}

\begin{milestone}[Step 6: richer PeTTa/MeTTa reasoning]
After the file-backed store is stable, load selected clusters into a PeTTa/MeTTa
reasoning context for consistency checks, simple inference, and later PLN inference
over the memory AtomSpace.

Tasks:
\begin{enumerate}
  \item Load active question and commitment atoms into a temporary AtomSpace.
  \item Add rules for detecting contradictions, e.g. active and resolved status on
        the same id.
  \item Add rules for surfacing overdue or blocked commitments.
  \item Add rules linking episodes, artifacts, and formalization notes.
  \item Add PLN-oriented representations for uncertain interpretive atoms such as
        hypotheses, importance estimates, confidence, and evidence relations.
\end{enumerate}

Exit criteria:
\begin{itemize}
  \item A test detects inconsistent statuses.
  \item A test retrieves all active commitments about a target project.
  \item A test runs at least one PLN-style inference over loaded medium-memory atoms.
  \item The reasoning layer remains optional; basic query still works without it.
\end{itemize}
\end{milestone}

\section{Validation plan}

\subsection{Unit tests}

Minimum tests:
\begin{enumerate}
  \item Append valid atom cluster.
  \item Reject malformed atom.
  \item Reject oversized atom.
  \item Query by id.
  \item Query by type.
  \item Query by about-entity.
  \item Query by status.
  \item Generate bounded prompt context.
\end{enumerate}

\subsection{Integration tests}

Use mock channels and mock providers first:
\begin{enumerate}
  \item Simulate a Ben request.
  \item Confirm an episode and commitment cluster are added.
  \item Simulate completion.
  \item Confirm the commitment status changes to resolved and links an artifact.
  \item Confirm prompt context includes the active item before completion and omits
        or compresses it after completion.
\end{enumerate}

\subsection{Live supervised smoke}

Only after mock tests pass:
\begin{enumerate}
  \item Run ProtomegaTron in supervised Telegram mode.
  \item Send one explicit memory-worthy request.
  \item Confirm the reply is concise and the medium-memory atom cluster is correct.
  \item Confirm no repeated idle backend calls occur.
  \item Stop or leave running only under the already-approved supervision policy.
\end{enumerate}

\section{Failure modes and mitigations}

\begin{risk}[Schema sprawl]
Too many predicates make memory less useful.
Mitigation: keep version 1 to episode, entity, commitment, question, hypothesis,
boundary, artifact, provenance, status, and about-links.
\end{risk}

\begin{risk}[False formality]
Bad atoms can look authoritative.
Mitigation: require provenance and confidence/status fields for interpretive atoms.
\end{risk}

\begin{risk}[Prompt bloat]
A growing medium memory could overwhelm context.
Mitigation: bounded prompt views, not full-store injection.
\end{risk}

\begin{risk}[Write amplification]
If every turn writes atoms, the store becomes noisy.
Mitigation: selective atomization guidance and tests that routine chatter writes nothing.
\end{risk}

\begin{risk}[Stale commitments]
Resolved or obsolete atoms may keep resurfacing.
Mitigation: explicit statuses and distillation markers.
\end{risk}

\section{Recommended first engineering patch}

The smallest useful patch is:
\begin{enumerate}
  \item Add \texttt{src/medium\_memory.py} with append, tail, and simple query.
  \item Add \texttt{src/medium\_memory.metta} exposing \texttt{medium-remember},
        \texttt{medium-query}, and \texttt{medium-active}.
  \item Add a config flag \texttt{enableMediumMemory} defaulting to false.
  \item Add tests using a temporary medium-memory file.
  \item Enable the feature only in the ProtomegaTron local wrapper after tests pass.
\end{enumerate}

This patch creates value without requiring a new database, new credentials, a new
service, or a global change in OmegaClaw behavior.

\section{Open design questions}

\begin{enumerate}
  \item Should \texttt{MM} be upstreamed into OmegaClaw Core after the local prototype,
        or remain ProtomegaTron-specific until the ontology stabilizes?
  \item Should the first query implementation remain text-based, or should it load the
        atom log into PeTTa on every query despite possible latency?
  \item What is the right decay policy: age-based, status-based, attention-based, or
        explicit distillation only?
  \item Should Hyperseed formalization notes be generated from \texttt{MM} atoms as a
        semi-automatic paper-writing substrate?
\end{enumerate}

\end{document}
